%content of the chapter classical electrodynamics
\chapter{Classical Electrodynamics}
\label{class-ed}
\section{Microscopic Maxwell Equations}
\label{micro-maxwell-eqs}

Omitting the long history of misconceptions and discoveries, 
one may say that, in addition to mass, particles may possess 
electric charge. Like mass, it is a scalar quantity, but unlike mass, 
charge may be positive, negative, or zero.
To describe the interaction of charged particles, 
it is necessary to introduce a new entity --- the electromagnetic field $ (\vec E(\vec r,t),\vec B(\vec r,t)) $.

The motion of  particle with charge $ e $ is described by Lorentz force
\begin{equation}
   \label {lorentz-force}
   \frac{d \vec p}{d t}=e \vec E(\vec r,t) + \frac{e}{c}\vec v \times \vec B(\vec r,t)
\end{equation}
Dynamics of electromagnetic field is described   Maxwell’s equations:
\begin{align}
\nabla \times \vec{E} + \frac{1}{c}\frac{\partial \vec{B}}{\partial t} &= 0, \label{me-faraday} \\
\nabla \cdot \vec{B} &= 0, \label{me-monopole} \\
\nabla \cdot \vec{E} &= 4\pi \rho, \label{me-gauss} \\
\nabla \times \vec{B} - \frac{1}{c}\frac{\partial \vec{E}}{\partial t} &= \frac{4\pi}{c}\vec{j}. \label{me-ampere}
\end{align}
Here 
\begin{equation}
   \label {dens1}
    \rho(\vec r,t)=\Sigma_a e_a \delta (\vec r-\vec r_a)
\end{equation}
and
\begin{equation}
   \label {curr1}
    \vec j(\vec r,t)=\Sigma_a e_a \vec v_a \delta (\vec r-\vec r_a)
\end{equation}
It is important to know that $ \delta $-function is actually not a function, but distribution
--- linear functional on the space of test functions. It is very useful in linear problems
but makes a lot of troubles in non-linear ones.

Taking the divergence of \eqref{me-ampere} and using \eqref{me-gauss} yields the continuity equation:
\begin{equation}
    \label{cont}
\partial_t \rho + \nabla \cdot \vec{j} = 0,
\end{equation}
demonstrating that charge conservation is not an independent 
postulate but a structural consequence of Maxwell’s equations.
In particular \eqref{dens1} and \eqref{curr1} satisfy \eqref{cont}.

\section{Maxwell Equations in Media}
\label{maxwell-eqs}


In media we can somehow divide charges and currents into two groups: 
bound ones which belong to media and free ones.
Bound charges can be expressed as
\begin{equation}
   \label {bound-charge}
   \rho_b=-\nabla \vec P
\end{equation}
where $ \vec P $  (polarization) is the dipole moment of unit volume
\begin{equation}
   \label {bound-current}
   \vec j_b =\frac{\partial \vec P}{\partial t}+ c \nabla \times \vec M
\end{equation}
where $ \vec M$ (magnetization) is the magnetic dipole moment of unit volume.

Thus macroscopic (averaged) charge and current densities become
\begin{equation}
   \label {tot-cur-ch}
   \vec j=\vec j_f+ \vec j_b \qquad \rho=\rho_f+\rho_b
\end{equation}
and macroscopic Maxwell equations read:
\begin{align}
    \nabla \cdot \vec{E} &= 4\pi (\rho_f -\nabla \vec P), \label{me-g1} \\
    \nabla \times \vec{B} - \frac{1}{c}\frac{\partial \vec{E}}{\partial t} &= 
    \frac{4\pi}{c}(\vec j_f+\frac{\partial \vec P}{\partial t}+ c \nabla \times \vec M). \label{me-a1}
\end{align}
or:
\begin{align}
    \nabla \cdot \vec D &= 4\pi \rho_f , \label{me-g2} \\
    \nabla \times \vec H - \frac{1}{c}\frac{\partial \vec D}{\partial t} &= \frac{4\pi}{c}\vec j_f . \label{me-a2}
\end{align}
Here $ \vec D = \vec E+4 \pi \vec P $ and $ \vec H = \vec E-4 \pi \vec M $.
In equations \eqref{me-g2} and \eqref{me-a2} $ \rho $ and $ \vec j $ are supposed to be smooth
functions as a result of averaging of \eqref{dens1} and \eqref{curr1}.
Homogeneous equations \eqref{me-faraday} and \eqref{me-monopole} remain the same.  

Some information about media should be added to make Maxwell
equations solvable. Usually it is done by introducing constitutive
relations and generalized Ohm's law which define somehow 
$ \vec H $ , $ \vec D $ and $ \vec j $ in terms of 
$ \vec E $ and $ \vec B $. A lot of interesting physics
is contained in specific form of these relations.

\section{Gauge Transformations}
\label{gauge-transformations}

It is convenient to introduce scalar $\phi$ and vector $\vec A$ such that
\begin{equation}
   \label {potA}
   \vec B = \nabla \times \vec A \quad \vec E = -\frac{1}{c} \frac{\partial \vec A}{\partial t} - \nabla \phi
\end{equation}
so the equations \eqref{me-monopole} and \eqref{me-faraday} become identities, remaining equations will
take the form
\begin{equation}
   \label {me-pot}
   \begin{split}
       \Box \vec A + \nabla ( \nabla \cdot \vec A+ \frac{1}{c} \dot \phi) &= \frac{4 \pi}{c} \vec j \\
       \Box \phi -\frac{1}{c} \frac{\partial}{\partial t}( \nabla \cdot \vec A+ \frac{1}{c} \dot \phi) & = 4 \pi \rho
   \end{split}
\end{equation}
Here $ \Box $ is d'alambert operator defined as \[  \frac{1}{c^2} \frac{\partial^2}{\partial t^2}-\nabla^2 \]

Equations \eqref{potA} do not define potentials uniquely, gauge transformation: 
\begin{equation}
   \label {gauge-transformation}
   \begin{split}
       \vec A &\rightarrow \vec A+ \nabla \chi \\
       \phi & \rightarrow \phi - \dot \chi /c
   \end{split}
\end{equation}
 where $ \chi  $ is arbitrary leaves fields $ \vec B $
and $ \vec E $ unchanged. So we can prescribe additional
condition (gauge) on potentials.

If we prescribe the condition :
\begin{equation}
   \label {lg}
    \nabla \cdot \vec A+ \frac{1}{c} \dot \phi=0
\end{equation}
which is evidently compatible with \eqref{cont} equations \eqref{me-pot} separate and read:
\begin{equation}
   \label {me-pot-lg}
   \begin{split}
       \Box \vec A &= -\frac{4 \pi}{c} \vec j \\
       \Box \phi  & = -4 \pi \rho
   \end{split}
\end{equation}


The Lorenz gauge \eqref{lg} is the most popular one, it
allows us to uncouple equations \eqref{me-pot-lg}   so that the scalar and vector
potentials are described by symmetrical uncoupled wave equations,
which is a characteristic of this gauge. 
%It is the only covariant gauge.
However even Lorenz condition leaves some arbitrariness: $ \chi \rightarrow \chi + \psi $ 
 in \eqref{gauge-transformation} changes nothing if $ \Box \psi=0 $.

A less popular gauge in textbooks is the Coulomb gauge.
\begin{equation}
   \label {cg}
   \nabla \cdot \vec A=0
\end{equation}
In this gauge the scalar potential satisfies an instantaneous
Poisson equation, which is a peculiar characteristic of this
gauge.
A number of other gauges was studied in \cite{heras}:
Kirchhoff gauge:
\begin{equation}
   \label {kg}
    \nabla \cdot \vec A- \frac{1}{c} \dot \phi=0
\end{equation}
which can be seen as Lorenz gauge in Euclidean 4-space.
Velocity gauge:
\begin{equation}
   \label {vg}
    \nabla \cdot \vec A- \frac{1}{v} \dot \phi=0
\end{equation}
The velocity gauge is one in which the scalar
potential propagates with an arbitrary speed. This gauge is
not well known despite the fact that it was proposed several
years ago. The v-gauge is really a family of gauges that
contains the Lorenz and Coulomb gauges as particular cases
and also includes the Kirchhoff gauge.

Temporal gauge:
\begin{equation}
   \label {tg}
   \phi=0
\end{equation}
 The temporal gauge is one in which 
the scalar potential is identically zero, which means that 
the electric and magnetic fields are defined 
only by the vector potential.

It is shown in \cite{heras} that  all these  potentials yield the 
same  retarded
electric and magnetic fields. 

\section {Covariant Form of Maxwell Equations}
\label{me-cov}

Introducing (formally) fourth coordinate $ x^0=c t $ we can define
 Minkowski space as  the manifold $M \cong \mathbb{R}^4$ 
equipped with the  metric
\begin{equation}
   \label {minkg}
    g_{i j} = g^{i j}=\mathrm{diag}(1,-1,-1,-1).
\end{equation} 
The meaning of it will be discussed later, meanwhile
it is used only to simplify notations. So
\begin{align}
    x^i&=(ct,\vec x) \label{x4} \\
    A^i&=(\phi,\vec A) \label{a4}\\
    j^i&=(c \rho, \vec j) \label{j4}\\
    \partial_i&= \frac{\partial}{\partial x^i} \label{d4}
\end{align}
So $ \Box=g^{i k} \partial_i \partial_k $ 
Lorenz gauge \eqref{lg} reads:
\begin{equation}
   \label {lgc}
   \partial_i A^i=0
\end{equation} 
 continuity equation \eqref{cont} becomes
\begin{equation}
   \label {contc}
   \partial_i j^i=0
\end{equation} 
and Maxwell equations \eqref{me-pot} take the form:
\begin{equation}
   \label {me-potc}
    \Box A^i= \frac{4 \pi}{c} j^i
\end{equation}

The fields $ (\vec E,\vec B) $ in covariant formalism
become parts of antisymmetric tensor of second rank:
\begin{equation}
   \label {ebc}
    F_{i k}=\partial_i A_k-\partial_k A_i
\end{equation}
where $ A_k=g_{k i} A^i=(\phi,-\vec A) $ so equations \eqref{me-gauss}
and \eqref{me-ampere} take the form
\begin{equation}
   \label {mec}
    \partial_k F^{i k}=-\frac{4 \pi}{c} j^k
\end{equation}

One can easily see that Maxwell equations save their form
under linear transformations preserving metric \eqref{minkg}.
However Galilean transformations change them, in particular
the velocity of light does not depend on inertial frame of
reference, which contradicts the rule of adding velocities.


It is often stated that the fields $ (\vec E,\vec B) $ or
$ F_{i k} $ are the physical variables and potentials
$A^i$ have no physical meaning. This statement based
on fact that potentials $ A^i $ are not gauge invariant.
However with argument is not sound, following this logic
we must declare canonical momenta as not physical since
it  is also not gauge invariant (see section \ref{gauge-invariance}).
In both cases there is specific gauge which gives direct physical
meaning, it is Lorenz gauge in case of potentials and
zero total derivative in Lagrangian in case of momenta.

\section{Lagrangian for Electromagnetic Field}
\label{lag-em}

Electromagnetic field is an example of continuous system, where
dynamic variables depend not only on $ t $ but also on $ \vec x $.
Such systems may be seen as infinite set of "usual" variables, so
that dependence on $ \vec x $ is a generalisation of index, numbering
"particles". Another (and I think more reasonable) approach is to
consider any component of the  field as as degree of freedom depending
on multidimensional "time".

Anyway Lagrange function is represented as space integral of 
Lagrange density and action becomes 4-dimensional integral.

Equations \eqref{mec} follow from action integral
\begin{equation}
   \label {em-ac-1}
    S_1=-\frac{1}{c^2} \int A_i j^i d^4 x-\frac{1}{16 \pi c} \int F_{ik} F^{ik} d^4 x
\end{equation}
and equations \eqref{me-potc} can be obtained from
action
Equations \eqref{mec} follow from action integral
\begin{equation}
   \label {em-ac-2}
    S_2=-\frac{1}{c^2} \int A_i j^i d^4 x-\frac{1}{8 \pi c} \int \partial_i A^k \partial^i A_k d^4 x 
\end{equation}
here potential $ A^i $ should be in Lorenz gauge \eqref{lgc}.

\section{Physics of Electromagnetic Field}
\label{phys-em}


Any solution of wave equation may be represented as linear combination of specific
solution of inhomogeneous equation and general solution of homogeneous one:
\begin{equation}
   \label {we-sol}
    A_i(x)=\int G_{ik}(x,y)j^k(y) d^4 y +A_i^{free} (x)
\end{equation}
where $  G_{ik}(x,y)$ is a Green's function.
On the other hand potential determines the motion of charges by Lorentz force \eqref{lorentz-force}
and hence currents. Generally the self-consistent problem can not be solved,
it can not be even correctly posed. 





